%============================================================
% CH1 · LESSON 1 — WHAT AN ARRAY IS
%============================================================
\subsection{What an Array Is: Memory, Indexing and Growth}
\label{ssec:ch1-array}

Every chapter of this book manipulates the same object: a row of values with positions $0, 1, \dots, n-1$. Before we learn to scan, compact and rotate it, we need to know what it \emph{is} inside the machine, because three facts about its memory layout explain most of what is fast and what is slow in the rest of the book:
\begin{enumerate}[leftmargin=*]
    \item reading or writing \texttt{a[i]} costs the same for every $i$ --- one address computation;
    \item walking through the array \emph{in order} is many times faster than jumping around in it;
    \item appending to a \texttt{std::vector} is cheap \emph{on average}, even though a single append can copy the whole array.
\end{enumerate}

%------------------------------------------------------------
\subsubsection*{1. Contiguous memory and address arithmetic}

\begin{definitionbox}[Array]
An \textbf{array} of length $n$ is a block of $n$ elements of the \emph{same type}, stored \emph{contiguously} (one directly after the other) in memory. If the first element starts at address \texttt{base} and one element occupies $s$ bytes, then
\[
\operatorname{address}(\texttt{a[i]}) \;=\; \texttt{base} + i \cdot s \qquad (0 \le i < n).
\]
\end{definitionbox}

The formula is the whole secret of $\bigO{1}$ indexing: the processor does not search for element $i$, it \emph{computes} where it is with one multiplication and one addition, whatever $i$ and $n$ are. This is also why indices start at 0 in C and C++: \texttt{a[i]} is literally \texttt{*(a + i)}, ``the element $i$ steps after the first'', and the first element is $0$ steps after itself.

\begin{diagrambox}
\begin{tikzpicture}
    \dpindex{ai}{1.35}{$i$}{0}{5}
    \dparray{a}{0.6}{\texttt{a}}{7/acell, 3/acell, 9/acell, 4/aread, 1/acell, 8/acell}
    \foreach \off [count=\i from 0] in {0,4,8,12,16,20}{
        \node[font=\scriptsize\ttfamily, text=darkgray] at (0.45+\i*0.9,-0.05) {+\off};
    }
    \node[rowlabel, font=\scriptsize\color{darkgray}] at (-0.15,-0.05) {bytes};
    \node[note, anchor=west, text width=5.6cm, align=left] at (5.8,0.6)
        {\texttt{int} has $s = 4$ bytes, so\\
         \texttt{a[3]} lives at $\texttt{base} + 3\cdot 4 = \texttt{base} + 12$.\\[3pt]
         Same cost for \texttt{a[0]} and \texttt{a[999999]}: one multiply, one add.};
\end{tikzpicture}
\figcap{Six \texttt{int}s in one contiguous block. The offsets come from the program below.}
\end{diagrambox}

The following program prints those offsets. It also shows that the element size depends on the type: \texttt{a2[2]} of a \texttt{long long} array is $2\cdot 8 = 16$ bytes after \texttt{a2[0]}.

\begin{lstlisting}[style=cpp]
#include <bits/stdc++.h>
using namespace std;

int main() {
    int a[6] = {7, 3, 9, 4, 1, 8};
    long long a2[3] = {1, 2, 3};
    char* base = (char*)&a[0];
    for (int i = 0; i < 6; i++)          // prints offsets 0, 4, 8, 12, 16, 20
        cout << "a[" << i << "] = " << a[i] << " at base + " << (char*)&a[i] - base << " bytes\n";
    cout << "sizeof(int) = " << sizeof(int) << ", sizeof(long long) = " << sizeof(long long) << "\n";
    cout << "a2[2] at base + " << (char*)&a2[2] - (char*)&a2[0] << " bytes\n";   // 16
    cout << "*(a + 3) = " << *(a + 3) << "\n";                                  // 4, same as a[3]
}
\end{lstlisting}

\textbf{The price of contiguity.} Because the elements must stay side by side, an array cannot grow in place (the memory after it may belong to something else), and inserting or deleting in the middle shifts every element behind the position: $\bigO{n}$. A \emph{linked list} makes the opposite trade: insertion after a known node is $\bigO{1}$, but reaching element $i$ means following $i$ pointers.

\begin{diagrambox}
\begin{tikzpicture}[x=1cm, y=1cm,
    ln/.style={rectangle, draw=darkgray!70, fill=white, minimum width=8mm, minimum height=6mm, font=\small},
    lp/.style={rectangle, draw=darkgray!70, fill=lightgray, minimum width=4mm, minimum height=6mm}]
    \begin{scope}
        \dparray{b}{0}{array}{7/acell, 3/acell, 9/acell, 4/aread}
        \draw[dep] (b0.north) to[bend left=35] node[lbl, above]{one jump: base $+ 3s$} (b3.north);
    \end{scope}
    \begin{scope}[xshift=5.6cm]
        \node[ln] (n0) at (0,0) {7}; \node[lp, right=0pt of n0] (p0) {};
        \node[ln] (n1) at (1.9,0.7) {3}; \node[lp, right=0pt of n1] (p1) {};
        \node[ln] (n2) at (3.6,-0.5) {9}; \node[lp, right=0pt of n2] (p2) {};
        \node[ln, fill=currentfill, draw=accentorange] (n3) at (5.5,0.4) {4}; \node[lp, right=0pt of n3] (p3) {};
        \draw[flow] (p0.center) -- (n1.west); \draw[flow] (p1.center) -- (n2.west); \draw[flow] (p2.center) -- (n3.west);
        \node[note, anchor=north] at (2.9,-1.0) {linked list: element 3 needs 3 pointer hops, each to an unpredictable address};
    \end{scope}
\end{tikzpicture}
\figcap{Random access: an array computes the address, a linked list walks to it.}
\end{diagrambox}

%------------------------------------------------------------
\subsubsection*{2. Cache lines: why scanning in order is fast}

Main memory is slow compared with the processor: fetching a value from RAM takes on the order of a hundred nanoseconds, during which the processor could execute hundreds of instructions. The hardware hides this with \textbf{caches} --- small, fast memories close to the processor --- and it never fetches a single byte: it fetches a whole \textbf{cache line}, typically 64 bytes, i.e.\ 16 consecutive \texttt{int}s.

\begin{diagrambox}
\begin{tikzpicture}[x=0.55cm, y=1cm]
    \foreach \i in {0,...,23}{
        \pgfmathtruncatemacro{\line}{\i/8}
        \ifnum\line=0 \node[cell, minimum width=5.5mm, minimum height=6mm, fill=hitfill, font=\tiny] at (\i,0) {\i};
        \else\ifnum\line=1 \node[cell, minimum width=5.5mm, minimum height=6mm, fill=computedfill, font=\tiny] at (\i,0) {\i};
        \else \node[cell, minimum width=5.5mm, minimum height=6mm, fill=basefill, font=\tiny] at (\i,0) {\i};
        \fi\fi
    }
    \foreach \k/\c in {0/memoorange, 1/tabulationblue, 2/successgreen}{
        \draw[thick, \c, decorate, decoration={brace, mirror, amplitude=4pt}]
            (\k*8-0.45,-0.45) -- (\k*8+7.45,-0.45) node[midway, below=4pt, font=\scriptsize\bfseries, text=\c] {cache line \k};
    }
    \draw[flow] (-0.4,0.75) -- node[lbl, above]{sequential scan: 1 slow fetch, then 7 fast reads per line} (23.4,0.75);
\end{tikzpicture}
\figcap{A schematic with 8 elements per line (real lines hold 16 \texttt{int}s). Reading \texttt{a[0]} brings \texttt{a[1..7]} into the cache for free; the hardware even prefetches the next line when it sees a sequential pattern.}
\end{diagrambox}

Two consequences you will rely on for the rest of the book:
\begin{itemize}[leftmargin=*]
    \item A \textbf{sequential scan} of an array is about as fast as a computer can go: each fetched line serves 16 reads. A loop over $10^8$ \texttt{int}s runs in a fraction of a second.
    \item \textbf{Jumping around} (random indices, linked lists, or a 2D array traversed column by column) may pay a full memory trip per element.
\end{itemize}

The effect is easy to measure. The program below sums a $4096 \times 4096$ matrix (64 MB) twice: row by row, which visits consecutive addresses, and column by column, which jumps $4096 \cdot 4 = 16384$ bytes between consecutive reads.

\begin{lstlisting}[style=cpp]
#include <bits/stdc++.h>
using namespace std;

const int N = 4096;
static int g[N][N];                              // 64 MB, row-major

int main() {
    for (int i = 0; i < N; i++)
        for (int j = 0; j < N; j++) g[i][j] = (i ^ j) & 7;
    auto t0 = chrono::steady_clock::now();
    long long s1 = 0;
    for (int i = 0; i < N; i++)                  // row by row: consecutive addresses
        for (int j = 0; j < N; j++) s1 += g[i][j];
    auto t1 = chrono::steady_clock::now();
    long long s2 = 0;
    for (int j = 0; j < N; j++)                  // column by column: jumps of N*4 bytes
        for (int i = 0; i < N; i++) s2 += g[i][j];
    auto t2 = chrono::steady_clock::now();
    auto ms = [](chrono::steady_clock::duration d) {
        return (long long)chrono::duration_cast<chrono::milliseconds>(d).count(); };
    cout << "row order    " << ms(t1 - t0) << " ms (sum " << s1 << ")\n";
    cout << "column order " << ms(t2 - t1) << " ms (sum " << s2 << ")\n";
}
\end{lstlisting}

On the author's machine (\texttt{g++ -O2}) three runs gave $13$--$24$ ms for row order and $462$--$500$ ms for column order: the \emph{same} $16.7$ million additions, about $20$--$35$ times slower, only because of the access pattern. Both sums are $58720256$. Your numbers will differ; the ratio will not disappear.

\begin{insightbox}[The practical rule]
Prefer algorithms that read arrays \emph{left to right}. When a 2D array is involved, make the \emph{last} index the inner loop. When two algorithms have the same $\bigO{\cdot}$, the one with sequential access usually wins by a constant factor of several times.
\end{insightbox}

%------------------------------------------------------------
\subsubsection*{3. Static arrays and \texttt{std::vector}}

C++ gives you two kinds of arrays.

\begin{table}[H]
\centering\small
\renewcommand{\arraystretch}{1.2}
\begin{tabularx}{\linewidth}{@{}>{\raggedright\arraybackslash}p{2.9cm} >{\raggedright\arraybackslash}X >{\raggedright\arraybackslash}X@{}}
\toprule
 & \textbf{static array} \texttt{int a[N];} & \textbf{dynamic array} \texttt{vector<int> v(n);} \\
\midrule
size fixed when & compiling (\texttt{N} is a constant) & running (\texttt{n} read from input); can grow \\
where it lives & global: static memory; local: the \emph{stack} & the heap (the vector object itself holds three pointers) \\
initial values & global: all zero; local: \emph{garbage} & zero (or the value you pass: \texttt{v(n, -1)}) \\
large sizes & global is fine; a local \texttt{int a[1000000]} can overflow the stack (often 1--8 MB) & fine \\
knows its size & no (\texttt{sizeof} trick only in scope) & \texttt{v.size()} \\
bounds check & never & only with \texttt{v.at(i)} \\
\bottomrule
\end{tabularx}
\end{table}

In contests both are common: a global \texttt{int a[200005];} sized by the maximal $n$ plus a few spare cells, or a \texttt{vector<int> a(n)} sized exactly. This book uses \texttt{vector} almost everywhere: it carries its size, it can be returned from functions, and it is what LeetCode passes to your code.

\begin{warningbox}[Uninitialised local arrays]
\begin{lstlisting}[style=snippet]
int main() {
    int cnt[100];            // WRONG: local, contains garbage
    int cnt2[100] = {};      // RIGHT: all zero
    vector<int> cnt3(100);   // RIGHT: all zero
}
\end{lstlisting}
A local counting array without \texttt{= \{\}} works on your machine ``by luck'' and fails on the judge.
\end{warningbox}

%------------------------------------------------------------
\subsubsection*{4. How \texttt{push\_back} grows a vector: amortised doubling}

A vector keeps a buffer of \textbf{capacity} $c$ elements of which the first \textbf{size} $s \le c$ are in use. \texttt{push\_back(x)} writes \texttt{x} into slot $s$ if $s < c$ --- that is $\bigO{1}$. If the buffer is full, the vector must
\begin{enumerate}[leftmargin=*]
    \item allocate a new, \emph{larger} buffer (libstdc++ uses $2c$; other libraries use $1.5c$),
    \item copy (or move) all $s$ old elements into it,
    \item free the old buffer, then write \texttt{x}.
\end{enumerate}

\begin{diagrambox}
\begin{tikzpicture}[x=0.62cm, y=0.95cm,
    bc/.style={cell, minimum width=6mm, minimum height=5.5mm, font=\scriptsize},
    bused/.style={bc, fill=computedfill},
    bnew/.style={bc, fill=currentfill, draw=accentorange, line width=1pt},
    bfree/.style={bc, fill=white, draw=darkgray!40}]
    \foreach \row/\cap/\used/\copied [count=\k from 0] in {0/1/0/0, 1/2/1/1, 2/4/2/2, 3/8/4/4}{
        \node[anchor=east, font=\scriptsize\bfseries\color{deepblue}] at (-0.6,-\k) {cap \cap};
        \foreach \j in {1,...,\cap}{
            \pgfmathtruncatemacro{\jj}{\j-1}
            \ifnum\jj<\used \node[bused] at (\jj,-\k) {\j};
            \else\ifnum\jj=\used \node[bnew] at (\jj,-\k) {\j};
            \else \node[bfree] at (\jj,-\k) {};
            \fi\fi
        }
        \node[anchor=west, font=\scriptsize, text=darkgray] at (8.4,-\k) {push \#\the\numexpr\used+1\relax: copy \copied{} old element(s), write 1};
    }
    \node[note, anchor=west, text width=5cm, align=left] at (8.4,-4.1)
        {pushes 5--8 then fill the free slots of the cap-8 buffer without any copy; push 9 copies 8 and doubles again};
\end{tikzpicture}
\figcap{Blue: elements copied from the old buffer. Orange: the new element. White: spare capacity, filled later for free.}
\end{diagrambox}

A single \texttt{push\_back} can therefore cost $\bigO{n}$. But it is cheap \textbf{on average}:

\begin{theorembox}[Amortised cost of doubling]
Starting from an empty vector, any sequence of $n$ \texttt{push\_back} calls with capacity doubling performs fewer than $2n$ element copies in total, so the total work is $\bigO{n}$ and the \textbf{amortised} cost per call is $\bigO{1}$.

\textbf{Proof.} Copies happen only when the size reaches a power of two, $1, 2, 4, \dots, 2^k$, where $2^k < n$ is the last such power below $n$; reaching size $2^j$ copies $2^j$ elements. The total is
\[
1 + 2 + 4 + \dots + 2^k \;=\; 2^{k+1} - 1 \;<\; 2 \cdot 2^k \;<\; 2n.
\]
Adding the $n$ writes of the new elements gives fewer than $3n$ element operations. \hfill$\blacksquare$
\end{theorembox}

The same argument works for any growth factor $g > 1$ (the copies form a geometric series $\le n/(g-1)$); it fails for \emph{additive} growth: growing by a constant $10$ slots copies $10 + 20 + 30 + \dots \approx n^2/20$ elements --- quadratic. The program below logs every reallocation:

\begin{lstlisting}[style=cpp]
#include <bits/stdc++.h>
using namespace std;

int main() {
    vector<int> v;
    size_t lastCap = v.capacity();
    long long copies = 0;                       // elements moved by reallocations
    for (int i = 1; i <= 1000; i++) {
        size_t before = v.size();
        v.push_back(i);
        if (v.capacity() != lastCap) {          // a reallocation happened
            copies += before;
            cout << "size " << v.size() << ": capacity " << lastCap << " -> " << v.capacity() << "\n";
            lastCap = v.capacity();
        }
    }
    cout << "1000 push_backs moved " << copies << " elements in total\n";
}
\end{lstlisting}

With g++ the capacities go $0 \to 1 \to 2 \to 4 \to \dots \to 512 \to 1024$ (11 reallocations) and the total number of moved elements is $1 + 2 + \dots + 512 = 1023 < 2 \cdot 1000$, exactly as the theorem predicts.

\textbf{\texttt{reserve}.} If you know the final size, \texttt{v.reserve(n)} allocates the buffer once and every later \texttt{push\_back} is a plain write. It changes the \emph{capacity}, not the \emph{size}:

\begin{warningbox}[\texttt{reserve} is not \texttt{resize}]
\begin{lstlisting}[style=snippet]
vector<int> v;
v.reserve(n);  v[0] = 5;     // WRONG: size is still 0, v[0] is out of range
vector<int> w;
w.resize(n);   w[0] = 5;     // RIGHT: size n, all zero
vector<int> u;
u.reserve(n);  u.push_back(5); // RIGHT: no reallocation for the first n pushes
\end{lstlisting}
\end{warningbox}

\begin{warningbox}[Pointers and references die on reallocation]
\begin{lstlisting}[style=snippet]
int& first = v[0];
v.push_back(42);          // may move the whole buffer ...
cout << first;            // ... WRONG: 'first' may point into freed memory
\end{lstlisting}
The same holds for iterators and for \texttt{\&v[i]}. Take indices instead of references when the vector may still grow.
\end{warningbox}

%------------------------------------------------------------
\subsubsection*{5. Two-dimensional arrays}

A grid of $R$ rows and $C$ columns can be stored in three ways.

\textbf{(a) Static 2D array} \texttt{int g[R][C]}: one contiguous block in \textbf{row-major} order --- row 0, then row 1, \ldots{} --- so
\[
\operatorname{address}(\texttt{g[i][j]}) = \texttt{base} + (i \cdot C + j)\cdot s .
\]
\textbf{(b) Flattened vector} \texttt{vector<int> g(R * C)} with \texttt{g[i * C + j]}: exactly the same layout, but sized at run time. The index $\mathit{id} = i \cdot C + j$ is reversible: $i = \mathit{id} / C$, $j = \mathit{id} \bmod C$ (Pattern~4, Section~\ref{sec:ch1-p4}, uses this a lot).

\textbf{(c) Vector of vectors} \texttt{vector<vector<int>> g(R, vector<int>(C))}: every row is a separate heap block. Rows may have different lengths (``jagged'' arrays), and \texttt{g[i]} can be swapped or passed around, but each access first follows the row pointer and consecutive rows are not adjacent.

\begin{diagrambox}
\begin{tikzpicture}[x=1cm, y=1cm]
    \begin{scope}
        \node[font=\scriptsize\bfseries, text=deepblue] at (1.35,1.95) {$3 \times 3$ grid};
        \foreach \r in {0,1,2}{
            \foreach \c in {0,1,2}{
                \pgfmathtruncatemacro{\v}{\r*3+\c}
                \ifnum\r=1 \node[cell, fill=hitfill, minimum width=8mm] at (0.45+\c*0.9,1.2-\r*0.75) {\v};
                \else \node[cell, minimum width=8mm] at (0.45+\c*0.9,1.2-\r*0.75) {\v};
                \fi
            }
        }
    \end{scope}
    \draw[flow] (3.0,0.45) -- node[lbl, above]{row-major} (4.3,0.45);
    \begin{scope}[xshift=4.6cm, yshift=0.45cm]
        \foreach \v in {0,...,8}{
            \ifnum\v>2 \ifnum\v<6
                \node[cell, fill=hitfill, minimum width=8mm] (f\v) at (0.4+\v*0.85,0) {\v};
            \else \node[cell, minimum width=8mm] (f\v) at (0.4+\v*0.85,0) {\v}; \fi
            \else \node[cell, minimum width=8mm] (f\v) at (0.4+\v*0.85,0) {\v}; \fi
            \node[idx] at (0.4+\v*0.85,0.6) {\v};
        }
        \draw[thick, memoorange, decorate, decoration={brace, mirror, amplitude=4pt}]
            ([yshift=-1.5mm]f3.south west) -- ([yshift=-1.5mm]f5.south east)
            node[midway, below=4pt, font=\scriptsize\bfseries, text=memoorange] {row 1: ids $1\cdot 3 + j$};
    \end{scope}
\end{tikzpicture}
\figcap{A $3 \times 3$ grid (cell value = its flat index) and its row-major layout. Cell $(1,2)$ has id $1 \cdot 3 + 2 = 5$; id $7$ is row $7/3 = 2$, column $7 \bmod 3 = 1$.}
\end{diagrambox}

\begin{lstlisting}[style=cpp]
#include <bits/stdc++.h>
using namespace std;

int main() {
    int R = 3, C = 4;
    vector<vector<int>> g(R, vector<int>(C));     // (c) rows are separate blocks
    vector<int> flat(R * C);                      // (b) one block, id = i*C + j
    for (int i = 0; i < R; i++)
        for (int j = 0; j < C; j++) {             // last index innermost: sequential
            g[i][j] = 10 * i + j;
            flat[i * C + j] = 10 * i + j;
        }
    int id = 7;
    cout << g[1][2] << " " << flat[1 * C + 2] << " | id " << id << " -> ("
         << id / C << ", " << id % C << ") = " << flat[id] << "\n";
    // prints: 12 12 | id 7 -> (1, 3) = 13
}
\end{lstlisting}

\begin{insightbox}[Which 2D layout to use]
\texttt{vector<vector<int>>} is the readable default and what LeetCode passes for grids. Switch to a flat \texttt{vector<int>} when the grid is large ($10^6$ cells and more), when you need a single integer id per cell (Pattern~4, graph searches on grids), or when every row must be contiguous for speed.
\end{insightbox}

\begin{checklistbox}[Lesson summary]
\begin{itemize}[leftmargin=*]
    \item[$\square$] \texttt{a[i]} is at $\texttt{base} + i\cdot s$: indexing is $\bigO{1}$; insertion in the middle is $\bigO{n}$.
    \item[$\square$] Scan in order; inner loop over the last index of a 2D array.
    \item[$\square$] Local arrays are uninitialised and live on a small stack; globals are zero and can be large.
    \item[$\square$] $n$ \texttt{push\_back}s cost $\bigO{n}$ in total (doubling); \texttt{reserve} removes the reallocations; references into a growing vector are unsafe.
    \item[$\square$] Row-major: $\mathit{id} = i\cdot C + j$, $i = \mathit{id}/C$, $j = \mathit{id} \bmod C$.
\end{itemize}
\end{checklistbox}
